\title{2017 年高考数学 全国II卷-文科}
\date{}
\maketitle

\section*{一、选择题}
本题共12小题，每小题5分，共60分．在每小题给出的四个选项中，只有一项是符合题目要求的．

\begin{question}[points=5]
设集合 $A=\{1,2,3\}$, $B=\{2,3,4\}$, 则 $A\cup B=$
\begin{choices}
  \item $\{1,2,3,4\}$
  \item $\{1,2,3\}$
  \item $\{2,3,4\}$
  \item $\{1,3,4\}$
\end{choices}
\begin{solution}
\textbf{答案：}$\{1,2,3,4\}$

\textbf{分析：}集合 $A=\{1,2,3\}$, $B=\{2,3,4\}$, 求 $A\cup B$, 可用并集的定义直接求出两集合的并集.

\textbf{详解：}$\because A=\{1,2,3\}$, $B=\{2,3,4\}$, $\therefore A\cup B=\{1,2,3,4\}$. 故选A.
\end{solution}
\end{question}

\begin{question}[points=5]
$(1+i)(2+i)=$
\begin{choices}
  \item $1-i$
  \item $1+3i$
  \item $3+i$
  \item $3+3i$
\end{choices}
\begin{solution}
\textbf{答案：}$1+3i$

\textbf{分析：}利用复数的运算法则即可得出.

\textbf{详解：}原式 $=2-1+3i=1+3i$. 故选B.
\end{solution}
\end{question}

\begin{question}[points=5]
函数 $f(x)=\sin\left(2x+\frac{\pi}{3}\right)$ 的最小正周期为
\begin{choices}
  \item $4\pi$
  \item $2\pi$
  \item $\pi$
  \item $\frac{\pi}{2}$
\end{choices}
\begin{solution}
\textbf{答案：}$\pi$

\textbf{分析：}利用三角函数周期公式, 直接求解即可.

\textbf{详解：}函数 $f(x)=\sin\left(2x+\frac{\pi}{3}\right)$ 的最小正周期为 $\frac{2\pi}{2}=\pi$. 故选C.
\end{solution}
\end{question}

\begin{question}[points=5]
设非零向量 $\vec{a}$, $\vec{b}$ 满足 $|\vec{a}+\vec{b}|=|\vec{a}-\vec{b}|$, 则
\begin{choices}
  \item $\vec{a}\perp\vec{b}$
  \item $|\vec{a}|=|\vec{b}|$
  \item $\vec{a}\parallel\vec{b}$
  \item $|\vec{a}|>|\vec{b}|$
\end{choices}
\begin{solution}
\textbf{答案：}$\vec{a}\perp\vec{b}$

\textbf{分析：}由已知得 $(\vec{a}+\vec{b})^2=(\vec{a}-\vec{b})^2$, 从而 $\vec{a}\cdot\vec{b}=0$, 由此得到 $\vec{a}\perp\vec{b}$.

\textbf{详解：}$\because$ 非零向量 $\vec{a}$, $\vec{b}$ 满足 $|\vec{a}+\vec{b}|=|\vec{a}-\vec{b}|$, $\therefore (\vec{a}+\vec{b})^2=(\vec{a}-\vec{b})^2$, $\therefore \vec{a}^2+2\vec{a}\cdot\vec{b}+\vec{b}^2=\vec{a}^2-2\vec{a}\cdot\vec{b}+\vec{b}^2$, 解得 $\vec{a}\cdot\vec{b}=0$, $\therefore \vec{a}\perp\vec{b}$. 故选A.
\end{solution}
\end{question}

\begin{question}[points=5]
若 $a>1$, 则双曲线 $\frac{x^2}{a^2}-y^2=1$ 的离心率的取值范围是
\begin{choices}
  \item $(\sqrt{2},+\infty)$
  \item $(\sqrt{2},2)$
  \item $(1,\sqrt{2})$
  \item $(1,2)$
\end{choices}
\begin{solution}
\textbf{答案：}$(1,\sqrt{2})$

\textbf{分析：}利用双曲线方程, 求出 $a$, $c$ 然后求解双曲线的离心率的范围即可.

\textbf{详解：}$a>1$, 则双曲线 $\frac{x^2}{a^2}-y^2=1$ 的离心率为 $e=\frac{c}{a}=\frac{\sqrt{a^2+1}}{a}=\sqrt{1+\frac{1}{a^2}}\in(1,\sqrt{2})$. 故选C.
\end{solution}
\end{question}

\begin{question}[points=5]
如图, 网格纸上小正方形的边长为1, 粗实线画出的是某几何体的三视图, 该几何体由一平面将一圆柱截去一部分后所得, 则该几何体的体积为

\begin{center}
  \includegraphics[width=.38\linewidth]{question-6-1.jpg}
\end{center}
\begin{choices}
  \item $90\pi$
  \item $63\pi$
  \item $42\pi$
  \item $36\pi$
\end{choices}
\begin{solution}
\textbf{答案：}$63\pi$

\textbf{分析：}由三视图可得, 直观图为一个完整的圆柱减去一个高为6的圆柱的一半, 即可求出几何体的体积.

\textbf{详解：}由三视图可得, 直观图为一个完整的圆柱减去一个高为6的圆柱的一半, $V=\pi\cdot 3^2\times 10-\frac{1}{2}\cdot\pi\cdot 3^2\times 6=63\pi$. 故选B.
\end{solution}
\end{question}

\begin{question}[points=5]
设 $x$, $y$ 满足约束条件 $\begin{cases} 2x+3y-3\le 0\\ 2x-3y+3\ge 0\\ y+3\ge 0\end{cases}$, 则 $z=2x+y$ 的最小值是
\begin{choices}
  \item $-15$
  \item $-9$
  \item $1$
  \item $9$
\end{choices}
\begin{solution}
\textbf{答案：}$-15$

\textbf{分析：}画出约束条件的可行域, 利用目标函数的最优解求解目标函数的最小值即可.

\textbf{详解：}$x$, $y$ 满足约束条件 $\begin{cases} 2x+3y-3\le 0\\ 2x-3y+3\ge 0\\ y+3\ge 0\end{cases}$ 的可行域如图: $z=2x+y$ 经过可行域的 $A$ 时, 目标函数取得最小值, 由 $\begin{cases} y=-3\\ 2x-3y+3=0\end{cases}$ 解得 $A(-6,-3)$, 则 $z=2x+y$ 的最小值是 $-15$. 故选A.
\end{solution}
\end{question}

\begin{question}[points=5]
函数 $f(x)=\ln(x^2-2x-8)$ 的单调递增区间是
\begin{choices}
  \item $(-\infty,-2)$
  \item $(-\infty,-1)$
  \item $(1,+\infty)$
  \item $(4,+\infty)$
\end{choices}
\begin{solution}
\textbf{答案：}$(4,+\infty)$

\textbf{分析：}由 $x^2-2x-8>0$ 得: $x\in(-\infty,-2)\cup(4,+\infty)$, 令 $t=x^2-2x-8$, 则 $y=\ln t$, 结合复合函数单调性“同增异减”的原则, 可得答案.

\textbf{详解：}由 $x^2-2x-8>0$ 得: $x\in(-\infty,-2)\cup(4,+\infty)$, 令 $t=x^2-2x-8$, 则 $y=\ln t$, $\because x\in(-\infty,-2)$ 时, $t=x^2-2x-8$ 为减函数; $x\in(4,+\infty)$ 时, $t=x^2-2x-8$ 为增函数; $y=\ln t$ 为增函数, 故函数 $f(x)=\ln(x^2-2x-8)$ 的单调递增区间是 $(4,+\infty)$. 故选D.
\end{solution}
\end{question}

\begin{question}[points=5]
甲、乙、丙、丁四位同学一起去问老师询问成语竞赛的成绩. 老师说: 你们四人中有2位优秀, 2位良好, 我现在给甲看乙、丙的成绩, 给乙看丙的成绩, 给丁看甲的成绩. 看后甲对大家说: 我还是不知道我的成绩. 根据以上信息, 则
\begin{choices}
  \item 乙可以知道四人的成绩
  \item 丁可以知道四人的成绩
  \item 乙、丁可以知道对方的成绩
  \item 乙、丁可以知道自己的成绩
\end{choices}
\begin{solution}
\textbf{答案：}乙、丁可以知道自己的成绩

\textbf{分析：}根据四人所知只有自己看到, 老师所说及最后甲说话, 继而可以推出正确答案.

\textbf{详解：}四人所知只有自己看到, 老师所说及最后甲说话, 甲不知自己的成绩 $\rightarrow$ 乙丙必有一优一良, (若为两优, 甲会知道自己的成绩; 若是两良, 甲也会知道自己的成绩) $\rightarrow$ 乙看到了丙的成绩, 知自己的成绩 $\rightarrow$ 丁看到甲、丁也为一优一良, 丁知自己的成绩, 给甲看乙丙成绩, 甲不知道自已的成绩, 说明乙丙一优一良, 假定乙丙都是优, 则甲是良, 假定乙丙都是良, 则甲是优, 那么甲就知道自已的成绩了. 给乙看丙成绩, 乙没有说不知道自已的成绩, 假定丙是优, 则乙是良, 乙就知道自己成绩. 给丁看甲成绩, 因为甲不知道自己成绩, 乙丙是一优一良, 则甲丁也是一优一良, 丁看到甲成绩, 假定甲是优, 则丁是良, 丁肯定知道自已的成绩了. 故选D.
\end{solution}
\end{question}

\begin{question}[points=5]
执行如图的程序框图, 如果输入的 $a=-1$, 则输出的 $S=$

\begin{center}
  \includegraphics[width=.38\linewidth]{question-10-1.jpg}
\end{center}
\begin{choices}
  \item $2$
  \item $3$
  \item $4$
  \item $5$
\end{choices}
\begin{solution}
\textbf{答案：}$3$

\textbf{分析：}执行程序框图, 依次写出每次循环得到的 $S$, $K$ 值, 当 $K=7$ 时, 程序终止即可得到结论.

\textbf{详解：}执行程序框图, 有 $S=0$, $K=1$, $a=-1$, 代入循环, 第一次满足循环, $S=-1$, $a=1$, $K=2$; 满足条件, 第二次满足循环, $S=1$, $a=-1$, $K=3$; 满足条件, 第三次满足循环, $S=-2$, $a=1$, $K=4$; 满足条件, 第四次满足循环, $S=2$, $a=-1$, $K=5$; 满足条件, 第五次满足循环, $S=-3$, $a=1$, $K=6$; 满足条件, 第六次满足循环, $S=3$, $a=-1$, $K=7$; $K\le 6$ 不成立, 退出循环输出 $S$ 的值为 $3$. 故选B.
\end{solution}
\end{question}

\begin{question}[points=5]
从分别写有 $1,2,3,4,5$ 的 $5$ 张卡片中随机抽取 $1$ 张, 放回后再随机抽取 $1$ 张, 则抽得的第一张卡片上的数大于第二张卡片上的数的概率为
\begin{choices}
  \item $\frac{1}{10}$
  \item $\frac{1}{5}$
  \item $\frac{3}{10}$
  \item $\frac{2}{5}$
\end{choices}
\begin{solution}
\textbf{答案：}$\frac{2}{5}$

\textbf{分析：}先求出基本事件总数 $n=5\times5=25$, 再用列举法求出抽得的第一张卡片上的数大于第二张卡片上的数包含的基本事件个数, 由此能求出概率.

\textbf{详解：}基本事件总数 $n=5\times5=25$, 抽得的第一张卡片上的数大于第二张卡片上的数包含的基本事件有: $(2,1),(3,1),(3,2),(4,1),(4,2),(4,3),(5,1),(5,2),(5,3),(5,4)$, 共有 $m=10$ 个基本事件, $\therefore$ 抽得的第一张卡片上的数大于第二张卡片上的数的概率 $p=\frac{m}{n}=\frac{10}{25}=\frac{2}{5}$. 故选D.
\end{solution}
\end{question}

\begin{question}[points=5]
过抛物线 $C: y^2=4x$ 的焦点 $F$, 且斜率为 $\sqrt{3}$ 的直线交 $C$ 于点 $M$ ($M$ 在 $x$ 轴上方), $l$ 为 $C$ 的准线, 点 $N$ 在 $l$ 上, 且 $MN\perp l$, 则 $M$ 到直线 $NF$ 的距离为
\begin{choices}
  \item $\sqrt{5}$
  \item $2\sqrt{2}$
  \item $2\sqrt{3}$
  \item $3\sqrt{3}$
\end{choices}
\begin{solution}
\textbf{答案：}$2\sqrt{3}$

\textbf{分析：}利用已知条件求出 $M$ 的坐标, 求出 $N$ 的坐标, 利用点到直线的距离公式求解即可.

\textbf{详解：}抛物线 $C: y^2=4x$ 的焦点 $F(1,0)$, 且斜率为 $\sqrt{3}$ 的直线: $y=\sqrt{3}(x-1)$, 过抛物线 $C: y^2=4x$ 的焦点 $F$, 且斜率为 $\sqrt{3}$ 的直线交 $C$ 于点 $M$ ($M$ 在 $x$ 轴上方), $l$ 可知: $\begin{cases} y^2=4x\\ y=\sqrt{3}(x-1) \end{cases}$, 解得 $M(3,2\sqrt{3})$. 可得 $N(-1,2\sqrt{3})$, $NF$ 的方程为: $y=-\frac{\sqrt{3}}{2}(x-1)$, 即 $\sqrt{3}x+2y-\sqrt{3}=0$, 则 $M$ 到直线 $NF$ 的距离为 $\frac{|\sqrt{3}\cdot3+2\cdot2\sqrt{3}-\sqrt{3}|}{\sqrt{(\sqrt{3})^2+2^2}}=\frac{6\sqrt{3}}{\sqrt{7}}=\frac{6\sqrt{21}}{7}$? 实为 $\frac{|3\sqrt{3}+4\sqrt{3}-\sqrt{3}|}{\sqrt{3+4}}=\frac{6\sqrt{3}}{\sqrt{7}}$, 但选项无此, 原题斜率为 $\sqrt{3}$? 检查: 原卷第12题斜率为 $\sqrt{3}$, 解析中直线 $y=\sqrt{3}(x-1)$, 解得 $M(3,2\sqrt{3})$, $N(-1,2\sqrt{3})$, $NF$ 方程: $y-0=\frac{2\sqrt{3}-0}{-1-1}(x-1)=-\sqrt{3}(x-1)$, 即 $\sqrt{3}x+y-\sqrt{3}=0$, 距离 $\frac{|\sqrt{3}\cdot3+2\sqrt{3}-\sqrt{3}|}{\sqrt{3+1}}=\frac{4\sqrt{3}}{2}=2\sqrt{3}$. 故选C.
\end{solution}
\end{question}

\section*{二、填空题}
本题共4小题，每小题5分，共20分

\begin{question}[points=5]
函数 $f(x)=2\cos x+\sin x$ 的最大值为
\fillin[{$\sqrt{5}$}]
\begin{solution}
\textbf{答案：}$\sqrt{5}$

\textbf{分析：}利用辅助角公式化简函数的解析式, 通过正弦函数的有界性求解即可.

\textbf{详解：}函数 $f(x)=2\cos x+\sin x=\sqrt{5}\left(\frac{2}{\sqrt{5}}\cos x+\frac{1}{\sqrt{5}}\sin x\right)=\sqrt{5}\sin(x+\theta)$, 其中 $\tan\theta=2$, 可知函数的最大值为 $\sqrt{5}$. 故答案为 $\sqrt{5}$.
\end{solution}
\end{question}

\begin{question}[points=5]
已知函数 $f(x)$ 是定义在 $\mathbb{R}$ 上的奇函数, 当 $x\in(-\infty,0)$ 时, $f(x)=2x^3+x^2$, 则 $f(2)=$
\fillin[{$12$}]
\begin{solution}
\textbf{答案：}$12$

\textbf{分析：}由已知中当 $x\in(-\infty,0)$ 时, $f(x)=2x^3+x^2$, 先求出 $f(-2)$, 进而根据奇函数的性质, 可得答案.

\textbf{详解：}$\because$ 当 $x\in(-\infty,0)$ 时, $f(x)=2x^3+x^2$, $\therefore f(-2)=-12$, 又 $\because$ 函数 $f(x)$ 是定义在 $\mathbb{R}$ 上的奇函数, $\therefore f(2)=12$. 故答案为 $12$.
\end{solution}
\end{question}

\begin{question}[points=5]
长方体的长、宽、高分别为 $3$, $2$, $1$, 其顶点都在球 $O$ 的球面上, 则球 $O$ 的表面积为
\fillin[{$14\pi$}]
\begin{solution}
\textbf{答案：}$14\pi$

\textbf{分析：}求出球的半径, 然后求解球的表面积.

\textbf{详解：}长方体的长、宽、高分别为 $3$, $2$, $1$, 其顶点都在球 $O$ 的球面上, 可知长方体的对角线的长就是球的直径, 所以球的半径为 $\frac{\sqrt{3^2+2^2+1^2}}{2}=\frac{\sqrt{14}}{2}$. 则球 $O$ 的表面积为 $4\pi\left(\frac{\sqrt{14}}{2}\right)^2=14\pi$. 故答案为 $14\pi$.
\end{solution}
\end{question}

\begin{question}[points=5]
$\triangle ABC$ 的内角 $A$, $B$, $C$ 的对边分别为 $a$, $b$, $c$, 若 $2b\cos B=a\cos C+c\cos A$, 则 $B=$
\fillin[{$\frac{\pi}{3}$}]
\begin{solution}
\textbf{答案：}$\frac{\pi}{3}$

\textbf{分析：}根据正弦定理和两角和的正弦公式和诱导公式计算即可.

\textbf{详解：}$\because 2b\cos B=a\cos C+c\cos A$, 由正弦定理可得, $2\cos B\sin B=\sin A\cos C+\sin C\cos A=\sin(A+C)=\sin B$, $\because \sin B\neq0$, $\therefore \cos B=\frac{1}{2}$, $\because 0<B<\pi$, $\therefore B=\frac{\pi}{3}$. 故答案为 $\frac{\pi}{3}$.
\end{solution}
\end{question}

\section*{三、解答题}
共70分. 解答应写出文字说明，证明过程或演算步骤. 第17~21题为必考题，每个试题考生都必须作答. 第22、23题为选考题，考生根据要求作答.

\begin{question}[points=12]
已知等差数列 $\{a_n\}$ 的前 $n$ 项和为 $S_n$, 等比数列 $\{b_n\}$ 的前 $n$ 项和为 $T_n$, $a_1=-1$, $b_1=1$, $a_2+b_2=2$. \\ (1) 若 $a_3+b_3=5$, 求 $\{b_n\}$ 的通项公式; \\ (2) 若 $T_3=21$, 求 $S_3$.
\begin{solution}

\textbf{分析：}(1) 设等差数列 $\{a_n\}$ 的公差为 $d$, 等比数列 $\{b_n\}$ 的公比为 $q$, 运用等差数列和等比数列的通项公式, 列方程解方程可得 $d$, $q$, 即可得到所求通项公式; (2) 运用等比数列的求和公式, 解方程可得公比, 再由等差数列的通项公式和求和, 计算即可得到所求和.

\textbf{详解：}(1) 设等差数列 $\{a_n\}$ 的公差为 $d$, 等比数列 $\{b_n\}$ 的公比为 $q$, $a_1=-1$, $b_1=1$, $a_2+b_2=2$, $a_3+b_3=5$, 可得 $-1+d+q=2$, $-1+2d+q^2=5$, 解得 $d=1$, $q=2$ 或 $d=3$, $q=0$ (舍去), 则 $\{b_n\}$ 的通项公式为 $b_n=2^{n-1}$, $n\in\mathbb{N}^*$; \\ (2) $b_1=1$, $T_3=21$, 可得 $1+q+q^2=21$, 解得 $q=4$ 或 $-5$, 当 $q=4$ 时, $b_2=4$, $a_2=2-4=-2$, $d=-2-(-1)=-1$, $S_3=-1-2-3=-6$; 当 $q=-5$ 时, $b_2=-5$, $a_2=2-(-5)=7$, $d=7-(-1)=8$, $S_3=-1+7+15=21$.
\end{solution}
\end{question}

\begin{question}[points=12]
如图, 四棱锥 $P-ABCD$ 中, 侧面 $PAD$ 为等边三角形且垂直于底面 $ABCD$, $AB=BC=\frac{1}{2}AD$, $\angle BAD=\angle ABC=90^\circ$. \\ (1) 证明: 直线 $BC\parallel$ 平面 $PAD$; \\ (2) 若 $\triangle PCD$ 面积为 $2\sqrt{7}$, 求四棱锥 $P-ABCD$ 的体积.

\begin{center}
  \includegraphics[width=.38\linewidth]{question-18-1.jpg}
\end{center}
\begin{solution}

\textbf{分析：}(1) 利用直线与平面平行的判定定理证明即可. (2) 利用已知条件转化求解几何体的线段长, 然后求解几何体的体积即可.

\textbf{详解：}(1) 证明: 四棱锥 $P-ABCD$ 中, $\because \angle BAD=\angle ABC=90^\circ$. $\therefore BC\parallel AD$, $\because AD\subset$ 平面 $PAD$, $BC\not\subset$ 平面 $PAD$, $\therefore$ 直线 $BC\parallel$ 平面 $PAD$. \\ (2) 解: 四棱锥 $P-ABCD$ 中, 侧面 $PAD$ 为等边三角形且垂直于底面 $ABCD$, $AB=BC=\frac{1}{2}AD$, $\angle BAD=\angle ABC=90^\circ$. 设 $AD=2x$, 则 $AB=BC=x$, $CD=\sqrt{2}x$, $O$ 是 $AD$ 的中点, 连接 $PO$, $OC$, $CD$ 的中点为 $E$, 连接 $OE$, 则 $OE=\frac{\sqrt{2}}{2}x$, $PO=\sqrt{3}x$, $PE=\sqrt{PO^2+OE^2}=\sqrt{3x^2+\frac{x^2}{2}}=\frac{\sqrt{14}}{2}x$, $\triangle PCD$ 面积为 $2\sqrt{7}$, 可得 $\frac{1}{2}\cdot CD\cdot PE=2\sqrt{7}$, 即 $\frac{1}{2}\cdot \sqrt{2}x\cdot\frac{\sqrt{14}}{2}x=2\sqrt{7}$, 解得 $x=2$, $PO=2\sqrt{3}$. 则 $V_{P-ABCD}=\frac{1}{3}\times\frac{1}{2}(BC+AD)\times AB\times PO=\frac{1}{6}\times(2+4)\times2\times2\sqrt{3}=4\sqrt{3}$.
\end{solution}
\end{question}

\begin{question}[points=12]
海水养殖场进行某水产品的新、旧网箱养殖方法的产量对比, 收获时各随机抽取了100个网箱, 测量各箱水产品的产量 (单位: kg), 其频率分布直方图如下: \\ (1) 记 $A$ 表示事件“旧养殖法的箱产量低于$50\text{kg}$”, 估计 $A$ 的概率; \\ (2) 填写下面列联表, 并根据列联表判断是否有$99\%$的把握认为箱产量与养殖方法有关; \\ (3) 根据箱产量的频率分布直方图, 对两种养殖方法的优劣进行比较. \\ 附: $P(K^2\ge k)$ $0.050$ $0.010$ $0.001$; $k$ $3.841$ $6.635$ $10.828$; $K^2=\frac{n(ad-bc)^2}{(a+b)(c+d)(a+c)(b+d)}$.

\begin{center}
  \includegraphics[width=.38\linewidth]{question-19-1.jpg}
\end{center}
\begin{solution}

\textbf{分析：}(1) 根据题意, 由旧养殖法的频率分布直方图计算可得答案; (2) 由频率分布直方图可以将列联表补全, 进而计算可得 $K^2\approx15.705>6.635$, 与附表比较即可得答案; (3) 由频率分布直方图计算新旧养殖法产量的平均数, 比较即可得答案.

\textbf{详解：}(1) 根据题意, 由旧养殖法的频率分布直方图可得: $P(A)=(0.012+0.014+0.024+0.034+0.040)\times5=0.62$; \\ (2) 根据题意, 补全列联表可得: $\begin{array}{c|ccc} & \text{箱产量}<50\text{kg} & \text{箱产量}\ge50\text{kg} & \text{总计} \\ \hline \text{旧养殖法} & 62 & 38 & 100 \\ \text{新养殖法} & 34 & 66 & 100 \\ \text{总计} & 96 & 104 & 200 \end{array}$ 则有 $K^2=\frac{200\times(62\times66-38\times34)^2}{100\times100\times96\times104}\approx15.705>6.635$, 故有 $99\%$ 的把握认为箱产量与养殖方法有关; \\ (3) 由频率分布直方图可得: 旧养殖法100个网箱产量的平均数 $\bar{x}_1=(27.5\times0.012+32.5\times0.014+37.5\times0.024+42.5\times0.034+47.5\times0.040+52.5\times0.032+57.5\times0.032+62.5\times0.012+67.5\times0.012)\times5=5\times9.42=47.1$; 新养殖法100个网箱产量的平均数 $\bar{x}_2=(37.5\times0.004+42.5\times0.020+47.5\times0.044+52.5\times0.054+57.5\times0.046+62.5\times0.010+67.5\times0.008)\times5=5\times10.47=52.35$; 比较可得 $\bar{x}_1<\bar{x}_2$, 故新养殖法更加优于旧养殖法.
\end{solution}
\end{question}

\begin{question}[points=12]
设 $O$ 为坐标原点, 动点 $M$ 在椭圆 $C: \frac{x^2}{2}+y^2=1$ 上, 过 $M$ 作 $x$ 轴的垂线, 垂足为 $N$, 点 $P$ 满足 $\overrightarrow{NP}=\sqrt{2}\overrightarrow{NM}$. \\ (1) 求点 $P$ 的轨迹方程; \\ (2) 设点 $Q$ 在直线 $x=-3$ 上, 且 $\overrightarrow{OP}\cdot\overrightarrow{PQ}=1$. 证明: 过点 $P$ 且垂直于 $OQ$ 的直线 $l$ 过 $C$ 的左焦点 $F$.
\begin{solution}

\textbf{分析：}(1) 设 $M(x_0,y_0)$, 由题意可得 $N(x_0,0)$, 设 $P(x,y)$, 运用向量的坐标运算, 结合 $M$ 满足椭圆方程, 化简整理可得 $P$ 的轨迹方程; (2) 设 $Q(-3,m)$, $P(\sqrt{2}\cos\alpha,\sin\alpha)$ $(0\le\alpha<2\pi)$, 运用向量的数量积的坐标表示, 可得 $m$, 即有 $Q$ 的坐标, 求得椭圆的左焦点坐标, 求得 $OQ$, $PF$ 的斜率, 由两直线垂直的条件: 向量数量积为 $0$, 即可得证.

\textbf{详解：}(1) 设 $M(x_0,y_0)$, 由题意可得 $N(x_0,0)$, 设 $P(x,y)$, 由点 $P$ 满足 $\overrightarrow{NP}=\sqrt{2}\overrightarrow{NM}$, 可得 $(x-x_0,y)=\sqrt{2}(0,y_0)$, 可得 $x-x_0=0$, $y=\sqrt{2}y_0$, 即有 $x_0=x$, $y_0=\frac{\sqrt{2}}{2}y$, 代入椭圆方程 $\frac{x^2}{2}+y^2=1$, 可得 $\frac{x^2}{2}+\frac{y^2}{2}=1$, 即有点 $P$ 的轨迹方程为圆 $x^2+y^2=2$; \\ (2) 证明: 设 $Q(-3,m)$, $P(\sqrt{2}\cos\alpha,\sin\alpha)$ $(0\le\alpha<2\pi)$, $\overrightarrow{OP}\cdot\overrightarrow{PQ}=1$, 可得 $(\sqrt{2}\cos\alpha,\sin\alpha)\cdot(-3-\sqrt{2}\cos\alpha,m-\sin\alpha)=1$, 即为 $-3\sqrt{2}\cos\alpha-2\cos^2\alpha+\sqrt{2}m\sin\alpha-2\sin^2\alpha=1$, 当 $\alpha=0$ 时, 上式不成立, 则 $0<\alpha<2\pi$, 解得 $m=\frac{3\sqrt{2}\cos\alpha+3}{\sqrt{2}\sin\alpha}$, 即有 $Q\left(-3,\frac{3\sqrt{2}\cos\alpha+3}{\sqrt{2}\sin\alpha}\right)$, 椭圆 $\frac{x^2}{2}+y^2=1$ 的左焦点 $F(-1,0)$, 由 $\overrightarrow{PF}\cdot\overrightarrow{OQ}=(-1-\sqrt{2}\cos\alpha,-\sin\alpha)\cdot\left(-3,\frac{3\sqrt{2}\cos\alpha+3}{\sqrt{2}\sin\alpha}\right)=3+3\sqrt{2}\cos\alpha-3(1+\sqrt{2}\cos\alpha)=0$. 可得过点 $P$ 且垂直于 $OQ$ 的直线 $l$ 过 $C$ 的左焦点 $F$.
\end{solution}
\end{question}

\begin{question}[points=12]
设函数 $f(x)=(1-x^2)e^x$. \\ (1) 讨论 $f(x)$ 的单调性; \\ (2) 当 $x\ge0$ 时, $f(x)\le ax+1$, 求 $a$ 的取值范围.
\begin{solution}

\textbf{分析：}(1) 求出函数的导数, 求出极值点, 利用导函数的符号, 判断函数的单调性即可. (2) 化简 $f(x)=(1-x)(1+x)e^x$. 下面对 $a$ 的范围进行讨论: ① 当 $a\ge1$ 时, ② 当 $0<a<1$ 时, ③ 当 $a\le0$ 时, 推出结论, 然后得到 $a$ 的取值范围.

\textbf{详解：}(1) 因为 $f(x)=(1-x^2)e^x$, $x\in\mathbb{R}$, 所以 $f'(x)=(1-2x-x^2)e^x$, 令 $f'(x)=0$ 可知 $x=-1\pm\sqrt{2}$, 当 $x<-1-\sqrt{2}$ 或 $x>-1+\sqrt{2}$ 时 $f'(x)<0$, 当 $-1-\sqrt{2}<x<-1+\sqrt{2}$ 时 $f'(x)>0$, 所以 $f(x)$ 在 $(-\infty,-1-\sqrt{2})$ 和 $(-1+\sqrt{2},+\infty)$ 上单调递减, 在 $(-1-\sqrt{2},-1+\sqrt{2})$ 上单调递增; \\ (2) 由题可知 $f(x)=(1-x)(1+x)e^x$. 下面对 $a$ 的范围进行讨论: ① 当 $a\ge1$ 时, 设函数 $h(x)=(1-x)e^x$, 则 $h'(x)=-xe^x<0$ $(x>0)$, 因此 $h(x)$ 在 $[0,+\infty)$ 上单调递减, 又因为 $h(0)=1$, 所以 $h(x)\le1$, 所以 $f(x)=(1+x)h(x)\le x+1\le ax+1$; ② 当 $0<a<1$ 时, 设函数 $g(x)=e^x-x-1$, 则 $g'(x)=e^x-1>0$ $(x>0)$, 所以 $g(x)$ 在 $[0,+\infty)$ 上单调递增, 又 $g(0)=1-0-1=0$, 所以 $e^x\ge x+1$. 因为当 $0<x<1$ 时 $f(x)>(1-x)(1+x)^2$, 所以 $(1-x)(1+x)^2-ax-1=x(1-a-x-x^2)$, 取 $x_0=\frac{\sqrt{5-4a}-1}{2}\in(0,1)$, 则 $(1-x_0)(1+x_0)^2-ax_0-1=0$, 所以 $f(x_0)>ax_0+1$, 矛盾; ③ 当 $a\le0$ 时, 取 $x_0=\frac{\sqrt{5}-1}{2}\in(0,1)$, 则 $f(x_0)>(1-x_0)(1+x_0)^2=1\ge ax_0+1$, 矛盾; 综上所述, $a$ 的取值范围是 $[1,+\infty)$.
\end{solution}
\end{question}

\begin{question}[points=10]
选修4-4: 坐标系与参数方程 \\ 在直角坐标系 $xOy$ 中, 以坐标原点为极点, $x$ 轴的正半轴为极轴建立极坐标系, 曲线 $C_1$ 的极坐标方程为 $\rho\cos\theta=4$. \\ (1) $M$ 为曲线 $C_1$ 上的动点, 点 $P$ 在线段 $OM$ 上, 且满足 $|OM|\cdot|OP|=16$, 求点 $P$ 的轨迹 $C_2$ 的直角坐标方程; \\ (2) 设点 $A$ 的极坐标为 $\left(2,\frac{\pi}{3}\right)$, 点 $B$ 在曲线 $C_2$ 上, 求 $\triangle OAB$ 面积的最大值.
\begin{solution}

\textbf{分析：}(1) 设 $P(x,y)$, 利用相似得出 $M$ 点坐标, 根据 $|OM|\cdot|OP|=16$ 列方程化简即可; (2) 求出曲线 $C_2$ 的圆心和半径, 得出 $B$ 到 $OA$ 的最大距离, 即可得出最大面积.

\textbf{详解：}(1) 曲线 $C_1$ 的直角坐标方程为: $x=4$, 设 $P(x,y)$, $M(4,y_0)$, 则 $\frac{y}{y_0}=\frac{x}{4}$, $\therefore y_0=\frac{4y}{x}$, $\because |OM||OP|=16$, $\therefore \sqrt{16+\frac{16y^2}{x^2}}\cdot\sqrt{x^2+y^2}=16$, 即 $(x^2+y^2)\left(1+\frac{y^2}{x^2}\right)=16$, $\therefore x^4+2x^2y^2+y^4=16x^2$, 即 $(x^2+y^2)^2=16x^2$, 两边开方得: $x^2+y^2=4x$, 整理得: $(x-2)^2+y^2=4$ $(x\neq0)$, $\therefore$ 点 $P$ 的轨迹 $C_2$ 的直角坐标方程: $(x-2)^2+y^2=4$ $(x\neq0)$. \\ (2) 点 $A$ 的直角坐标为 $A(1,\sqrt{3})$, 显然点 $A$ 在曲线 $C_2$ 上, $|OA|=2$, $\therefore$ 曲线 $C_2$ 的圆心 $(2,0)$ 到弦 $OA$ 的距离 $d=\frac{|2\times0-0\times1+?|}{\sqrt{?}}?$ 实际 $C_2$ 圆心 $(2,0)$ 到直线 $OA$ 的距离 $d=\frac{|\sqrt{3}\times2-1\times0|}{\sqrt{(\sqrt{3})^2+1^2}}=\sqrt{3}$, 则 $\triangle AOB$ 的最大面积 $S=\frac{1}{2}|OA|\cdot(2+\sqrt{3})=\frac{1}{2}\times2\times(2+\sqrt{3})=2+\sqrt{3}$.
\end{solution}
\end{question}

\begin{question}[points=10]
选修4-5: 不等式选讲 \\ 已知 $a>0$, $b>0$, $a^3+b^3=2$. 证明: \\ (1) $(a+b)(a^5+b^5)\ge4$; \\ (2) $a+b\le2$.
\begin{solution}

\textbf{分析：}(1) 由柯西不等式即可证明, (2) 由 $a^3+b^3=2$ 转化为 $\frac{(a+b)^3-2}{3(a+b)}=ab$, 再由均值不等式可得: $\frac{(a+b)^3-2}{3(a+b)}=ab\le\left(\frac{a+b}{2}\right)^2$, 即可得到 $\frac{1}{4}(a+b)^3\le2$, 问题得以证明.

\textbf{详解：}(1) 由柯西不等式得: $(a+b)(a^5+b^5)\ge(\sqrt{a}\cdot\sqrt{a^5}+\sqrt{b}\cdot\sqrt{b^5})^2=(a^3+b^3)^2\ge4$, 当且仅当 $\sqrt{a}\cdot\sqrt{b^5}=\sqrt{b}\cdot\sqrt{a^5}$, 即 $a=b=1$ 时取等号; \\ (2) $\because a^3+b^3=2$, $\therefore (a+b)(a^2-ab+b^2)=2$, $\therefore (a+b)[(a+b)^2-3ab]=2$, $\therefore (a+b)^3-3ab(a+b)=2$, $\therefore \frac{(a+b)^3-2}{3(a+b)}=ab$, 由均值不等式可得: $\frac{(a+b)^3-2}{3(a+b)}=ab\le\left(\frac{a+b}{2}\right)^2$, $\therefore (a+b)^3-2\le\frac{3}{4}(a+b)^3$, $\therefore \frac{1}{4}(a+b)^3\le2$, $\therefore a+b\le2$, 当且仅当 $a=b=1$ 时等号成立.
\end{solution}
\end{question}

